\subsection{The literal colored last-meridian comparison}

\begin{statement}{D23.P32}{Proved}{Three colored raw lifts remain over the single mixed normal form}
\end{statement}
Assume an actual normalized multiplier-one full-GT element has
$\rho\equiv0\pmod3$, $a_3=1$, and $a_2=0$.
After the source normalization of C34, suppose its original
Hurwitz data has label $(1,0)$. Its actual action on the
colored degree-$5376$ orbit has a fixed-product raw lift of
degree $32256$. Among the canonical linear lifts of that
projective action, exactly three pass the following necessary
conditions simultaneously:
\begin{enumerate}
\item commutation with the first half twist and the first
three-strand full twist;
\item F40's literal last-meridian formula, with the marked
cubic value of P10 and the eight identity values forced by C34;
\item the same-kernel value identifications, the two-cycle,
and the three-cycles for $m=0$ and $m=-1$ at every peripheral
orientation of each relevant marked kernel.
\end{enumerate}
The three possibilities have the sixteen canonical values
in Table~\ref{D23.P32.table}. Each of these sixteen-tuples
belongs to the simultaneous image of the derived free group.
For each, at least one complete extension to all thirty-eight
peripheral marked kernels also belongs to that derived image.
Thus this colored comparison sharpens the finite data attached
to the last normal form, but does not exclude that form.
No surviving finite value is asserted to lift to an actual
full-GT element.

\begin{proof}
\medskip\noindent\textbf{The raw colored orbit and its actual action.}\quad
Use P29's literal five-tuple orbit, with four entries in
$C_+$ and one in $C_-$, product one, modulo simultaneous
$E_0$ conjugacy. Its original sector has degree $1344$;
the other four sectors form the degree-$5376$ colored
$B_4$ orbit. The first four entries are exactly the
four-tuples used in the colored constructor. Their products
belong to the order-six class $C_-$.

Fix the representative
\[
 p=(0,1,2,0,1,2)\in E_0
\]
in the translation-then-matrix encoding
$(v_1,v_2,A_{11},A_{12},A_{21},A_{22})$ over $\mathbb F_3$.
Its centralizer is the cyclic group $\langle p\rangle$ of
order six. For each conjugacy class of tuples choose one
conjugate with product $p$, and then take its six conjugates
by $p^k$. The action on these tuples is free, since they
generate the center-free group $E_0$. This gives exactly
$6\cdot5376=32256$ distinct literal tuples. The three
positive Hurwitz generators satisfy the Artin relations,
project to the stored colored action, and have full twist
equal to the six-cycle deck transformation. In particular
the cover is connected.

At multiplier one, the actual tuple-kernel action fixes the
peripheral product exactly by F40. It commutes with
postcomposition by $E_0$, and hence with the deck action.
It preserves the raw preimage of the colored orbit: the
colored orbit is the unique orbit of its degree in the
complete union of the four color positions, and the argument
of P18 applies to this union. Alternatively its preservation
follows from the actual full spherical action in P29.
That spherical action has the four generator covariances of
P30 and its already known restriction on the original sector.
The connected labeled graph in P30 therefore determines it
uniquely. Its restriction to the colored sector is exactly
the permutation $V$ used here.

\medskip\noindent\textbf{Why the six linear lifts per block are necessary.}\quad
There are forty-eight projective blocks of degree $112$.
Fix one block and one exact tuple of linear parts
$A_1,\ldots,A_4\in SL_2(\mathbb F_3)$. Every $I-A_i$
and $I-A_1A_2A_3A_4$ is invertible. If
$g_i=(v_i,A_i)$, set
\[
 q_i=(I-A_i)^{-1}v_i,\qquad
 w=(q_1-q_4,q_2-q_4,q_3-q_4)\in\mathbb F_3^6.
\]
These are the same linear cohomology coordinates as in P23;
that argument uses invertibility, not equality of the four
peripheral classes. For completeness, a tuple of translations
is a crossed homomorphism of the free kernel into the given
linear module. Changing all $q_i$ by a common vector is a
coboundary. The condition that the product be the prescribed
$p$ selects one representative of each cohomology class,
because $I-A_1A_2A_3A_4$ is invertible. Thus $w$ gives
linear coordinates on the entire six-dimensional space.
The zero class is the tuple whose four fixed points equal
the unique fixed point of $p$.

Precomposition by the actual kernel automorphism induces a
linear map on this cohomology. Its new linear tuple depends
only on the old linear tuple, not on $w$. Since $V$ fixes
the projective block, the two linear tuples are simultaneously
$SL_2(\mathbb F_3)$-conjugate. The conjugating matrix commutes
with the linear part of $p$; its unique affine correction
fixing the fixed point of $p$ lies in $C_{E_0}(p)$.
After one constant deck shift, the chosen linear framing
is therefore preserved and the lift is linear on $w$.
This is the necessary linearity assertion; no arbitrary
point-dependent deck phase is permitted.

The exact raw inventory gives three linear framings per
projective block, each containing $224$ nonzero vectors,
with each projective point represented by the pair $w,-w$.
The deck cube is exactly $w\mapsto-w$. Choose six independent
vectors. A projective lift sends them to their prescribed
lines, so, after fixing the first sign, only $2^5=32$
matrices can occur. The certificate exhausts these matrices
and verifies the action on all $224$ vectors. Precisely one
passes. Its negative gives the other lift in that framing;
deck transport through the three framings gives exactly six
linear deck-equivariant lifts of the block action.

\medskip\noindent\textbf{Known values and the initial phase equations.}\quad
Let $a=p^{-1}$. P10 fixes the cubic marked value at
$\rho=0,a_3=1$ to be the translation $(1,1)$ with identity
linear part. C34 fixes the values at its eight marked
kernels to be the identity. The exact Cayley graph
identifications of C34 transport these nine values to all
pairs $(a,b)$ occurring here. The known kernels are
\[
                  37,42,49,58,61,66,73,77,78.
\]
They apply to $12384$ of the $32256$ literal tuples.
For a raw permutation $W$, the required equation is
\[
 \operatorname{last}(W^{-1}T)
       =f(a,\operatorname{last}(T))^{-1}
          \operatorname{last}(T)
          f(a,\operatorname{last}(T)).                 \tag{D23.P32.1}
\]
The inverse is essential because the left action is
precomposition by the inverse kernel automorphism.

Start with one canonical linear lift on each block and
write a general lift as its composition with a constant
deck shift $s_j\in\mathbb Z/6$ on block $j$.
Commutation with the first half twist and the first
three-strand full twist is an explicit linear potential
system over $\mathbb Z/6$. Both braid elements are fixed
by the actual multiplier-one action. The certificate checks
each pointwise edge equation, not just a selected point.
Together with (D23.P32.1) on the known kernels, the system
has twelve connected components, nine with a forced shift
and three with six choices. Thus exactly $6^3=216$
canonical linear raw lifts remain at this stage.

\medskip\noindent\textbf{All other last-entry values and all orientations.}\quad
The last entry ranges over $C_+\cup C_-$, of size $72$.
For each of the $216$ lifts first require that
$\operatorname{last}(W^{-1}T)$ depend only on
$\operatorname{last}(T)$. Exactly six lifts have this
necessary property. They are the assemblies numbered
$0,43,86,129,172,215$ in lexicographic order of the
three free shifts.

The pairs $(a,b)$ have sixteen marked kernels, namely the
sixteen rows in Table~\ref{D23.P32.table}. For each kernel
we exhaust its derived subgroup and retain exactly the
values producing the required conjugates in (D23.P32.1)
for all occurrences of that kernel. All six lifts have
nonempty fibers. Impose also the nine anchored values
already described.

Close these sixteen kernels under swapping and peripheral
three-cycling. There are thirteen orbits. The exhaustive
algorithm is the one proved in C34: for a prescribed
$g_0=f(x,y)$ enumerate every possible derived value
$g_1=f(y,(xy)^{-1})$, force
$g_2=(g_1g_0)^{-1}$, impose duality, transport all six
values by the exact marked isomorphisms, and check both
three-cycles at every canonical orientation. Values on a
prescribed kernel must belong to its entire allowed fiber,
so this step exhausts the fibers instead of selecting a
representative. The $m=-1$ equation is necessary because
$(1,f)(-1,1)=(-1,f)$ is an actual GT element with the
same associator word.

Exactly the assemblies $0,86,172$ pass. For each, restriction
of its possible thirteen orbit-extensions to the sixteen
original kernels is a singleton. These unique values are
shown in the table. An independent identity control, with
all associator values equal to one, passes every cycle
at every orientation.

\begin{table}[h]
\centering
\begin{tabular}{c|rrr}
canonical kernel&assembly $0$&assembly $86$&assembly $172$\\\hline
34&6&6&6\\
37&6&6&6\\
41&6&6&6\\
42&6&6&6\\
46&6&126&174\\
49&6&6&6\\
56&150&150&150\\
58&6&6&6\\
61&6&6&6\\
62&30&30&30\\
63&174&174&174\\
66&6&6&6\\
70&6&6&6\\
73&6&6&6\\
77&30&30&30\\
78&6&6&6
\end{tabular}
\caption{Canonical $E_0$ element indices in C34's exact
inventory; index $6$ is the identity. These are marked
word values, not unmarked conjugacy classes.}
\label{D23.P32.table}
\end{table}

\medskip\noindent\textbf{The complete sixteen-value common-word check.}\quad
Represent each canonical $E_0$ pair faithfully on the nine
affine points and take the disjoint union of the sixteen
marked actions. Its free-group image has degree $144$,
order $204073344$, and derived subgroup of order
$11337408$, hence derived index $18$.
Each of the three literal target permutations in the table
belongs to that derived subgroup. This proves simultaneous
ordinary commutator-word realizability of the sixteen values.

The full union of the thirteen peripheral orbits has
thirty-eight marked kernels. Its faithful disjoint affine
action has degree $342$, free-group image of order
$124937789194027008$, and derived subgroup of order
$3470494144278528$, of index $36$.
For each of the three assemblies, the possible extensions
already satisfying every displayed gentle equation number
$20736$. Testing these explicitly against the derived image
finds a witness after respectively $1$, $315$, and $158$
tests. The complete thirty-eight values of each witness are
retained in the certificate and their derived membership
is checked exactly. The search stops at the first witness;
no assertion that all $20736$ extensions are realizable is
needed or made. This proves the stronger common-word
assertion without producing a profinite associator.
Full raw colored braid covariance and correlations with
other source quotients are additional conditions.
\end{proof}

\medskip\noindent\textbf{Certificates and exact scope.}\quad
The complete construction is
\path{agents/pentagon/hurwitz_colored_raw_cover.py}.
The exhaustive $48$-block linear lift and phase equations are
\path{agents/pentagon/hurwitz_colored_raw_last_pair.py}.
The $216$-assembly and complete marked-orientation census is
\path{agents/pentagon/hurwitz_colored_last_gentle.py}.
The full sixteen-value simultaneous derived membership check is
\path{agents/pentagon/hurwitz_colored_last_commonword.py}.
The thirty-eight-kernel extension witnesses and exact derived
membership checks are
\path{agents/pentagon/hurwitz_colored_full_core.py}.
Each has JSON/NPZ records. P10 fixes the cubic marking;
C34 fixes the eight zero evaluations and proves the kernel
identification method; F40 fixes the literal free-factor
and inverse-action normalization; P18, P23 and P29--P30
supply the action and linear cohomology inputs used above.
The whole proof and finite certificates have independent review events recorded in the run ledger; the thirty-eight-kernel strengthening has also been independently read and replayed.
